\section{Conclusion}
\label{sec:conclusion}

When shard summaries go out of date, a split vector database can miss new matches.
The best use of a fixed network budget for finding them depends on where new data goes and its topics.

In our tested setting, the results suggest three rules of thumb.
First, when new data joins the most similar shard, spend the budget on extra search if shards have spare capacity.
Push is the other choice, if the coordinator can store and scan new passages.
Extra search won at every query rate on mixed topics and up to \RNumPlainQueriesDefault{} queries per step on new topics.
At \RNumPlainQueriesHigh{} queries on new topics, most-changed overtook it, but only fragilely once it paid for its scores.
Push won at every rate we tried, but at the default rate extra search raised overall recall far more.
Second, when each step's new data lands on one shard, refresh that shard if the data brings a new topic.
When topics are mixed, push instead, since fresh summaries there also lower overall recall.
When eight shards share the new data, push on mixed topics.
On a new topic, neither refresh nor push wins reliably at the smaller budgets, and push wins at the largest.
Third, to choose which shards to refresh, rank them by how much their data changed rather than by probing.
Our evidence comes from eight shards on new topics at budgets of four and eight, even with scores paid.
It also rests on recall alone, since we did not score most-changed with our ranking measure.

These results come from one simulator and one data set, with some costs uncharged.
Future work should test summaries that adapt to new data and coordinators that know where new data goes.
It should also test higher arrival rates and slow responses.
